\documentclass{article}

\title{Fault-Tolerant Long-Context LLM Serving via\\Checkpoint-Backed Preemption}
\author{Yi Zhong, Chen Yang, Guanghua Sun}
\date{April 29, 2026}

\begin{document}
\maketitle

\section*{Motivation}
Long-context retrieval-augmented generation (RAG) is increasingly common in production LLM serving, with prompts ranging from 32K to 128K tokens. In this regime, recovery is asymmetric: a mid-run engine failure forfeits 20--90 seconds of prefill compute on a single request, and naive re-prefill on a surviving replica cannot meet reasonable TTFT SLOs. Fault tolerance is a real production need.

Existing work treats fault tolerance and scheduling separately. Fault-tolerant LLM serving systems focus on recovery mechanisms (e.g., KV cache transfer in Mooncake~[6]) but assume static workloads. SLO-aware scheduling systems (QLM~[7], Scorpio~[8], JITServe~[9], SLOs-Serve~[10]) optimize admission and dispatch but assume no faults, and crucially cannot preempt long-context requests because reclaiming a mid-prefill request would lose tens of seconds of work. This separation hides an opportunity: the host-memory state maintained for fault recovery is exactly what an SLO-aware scheduler needs to preempt cheaply.

\section*{Research Question}
Can the host-memory KV state maintained for fault recovery serve as a \emph{dual-purpose primitive}---enabling both fast restore after failure and cheap preemption of in-flight long-context requests during normal operation? We study whether combining KV-checkpoint-based fault tolerance with preemption-aware SLO scheduling yields qualitatively different behavior from either alone, particularly under high post-fault utilization where a fast restore mechanism alone is insufficient because recovery requests still wait in the FCFS tail.

\section*{System Design}
We build on vLLM (v0.16) and add three components:
\begin{itemize}
    \item \textbf{KV cache checkpointing.} Periodic asynchronous copy of per-request KV blocks from GPU to host pinned memory, then publish to a shared memory file system (\texttt{/dev/shm}) using atomic write. Frequency is governed by an online cost model.
    \item \textbf{Cross-engine restore.} On engine failure, displaced requests are rerouted to a surviving engine, which restores KV blocks from the shared file system into freshly allocated GPU blocks and resumes decoding without re-prefilling the prompt.
    \item \textbf{Preemption-aware SLO scheduler.} At each scheduling step, the scheduler computes the remaining SLO budget (TTFT, TPOT) for every pending request and sorts by tightest budget first. A currently-executing request is preempted (its checkpointed KV state retained on host) when its budget significantly exceeds the top pending request, with a hysteresis margin to avoid oscillation. After a fault, all in-flight and recovered requests are immediately rescheduled.
\end{itemize}
The key design observation is that vanilla SLO-aware schedulers cannot preempt long-context requests cheaply because mid-prefill state is too expensive to abandon, while vanilla fault-tolerant systems already maintain that state but do not expose it to the scheduler. By making the checkpoint visible to both subsystems, preemption becomes a side-effect of work already paid for.

\section*{Evaluation Plan}
\textbf{Baselines.} We compare four configurations on the same vLLM substrate: \emph{No-FT + FCFS} (vanilla; in-flight requests lost on failure), \emph{Reprefill + FCFS} (reroute and re-prefill on the surviving engine), \emph{Checkpoint + FCFS} (KV checkpoint and restore with vanilla scheduling), and \emph{Checkpoint + SLO-aware (ours)} (checkpoint plus preemption-aware scheduling). The third vs.\ fourth comparison isolates the scheduling contribution; the fourth vs.\ second isolates the fault-tolerance contribution.

\textbf{Workloads.} ShareGPT for short-context chat traffic, LongBench~[12] (NarrativeQA, QMSum, MuSiQue) for long-context document understanding, and the Azure LLM inference trace~[5] for production-realistic arrival patterns. Each workload uses a Poisson arrival process with three random seeds for noise reduction.

\textbf{Faults.} A single mid-run engine failure is injected at 350~s into a 700~s cell, drawn from real production fault patterns. We also include a no-fault control to measure normal-path overhead.

\textbf{Metrics.} The primary metric is goodput, defined as SLO-met output tokens per second across the system. Secondary metrics include recovery success rate, TTFT and TPOT distributions, and system-wide GPU utilization. SLO settings are varied along TTFT and failover-gap axes.

\textbf{Hypothesis.} Under tight TTFT (below 2 seconds for 15K-token prompts), Checkpoint+FCFS will improve over Reprefill+FCFS but its mechanical advantage is partially hidden by queueing at the surviving engine. Checkpoint+SLO-aware preserves the advantage by allowing recovery requests to preempt non-urgent work. Under loose SLOs, all three FT configurations converge.

\section*{Related Work}
LLM serving work on scheduling focuses on steady-state SLO compliance: DistServe~[2], Sarathi-Serve~[3], and Llumnix~[4] introduce phase-aware dispatch, chunked prefill, and live migration respectively, but assume no faults. SLO-aware systems (QLM~[7], Scorpio~[8], JITServe~[9], SLOs-Serve~[10]) manage admission under load but do not preempt long-context requests because mid-prefill state cannot be cheaply reclaimed. Fault tolerance for LLM serving is less explored: Mooncake~[6] transfers KV cache across instances for prefix-aware routing, sharing mechanism with our restore path but targeting cache reuse; BanaServe~[11] rebalances placement under load imbalance, not faults; Splitwise~[5] releases the trace we use. Our work observes that the KV state maintained for fault tolerance is precisely what SLO-aware scheduling needs, and studies what becomes possible when both subsystems share this primitive.

\section*{Resources}
\textbf{Hardware.} 2+ GPUs configured for data-parallel (dp) deployment, sufficient to host multiple vLLM engine replicas and inject controlled engine failures across them.

\textbf{Software.} The serving stack is built on vLLM v0.16~[1] with the standard PyTorch and CUDA toolchain. Profiles for decode capacity and checkpoint cost are calibrated locally on the target hardware.

\textbf{Models and datasets.} We use the open-weight Llama-3.1-8B-Instruct model. Datasets include ShareGPT-Vicuna (5000 conversations), LongBench~[12], and the Azure LLM inference trace from the Splitwise release~[5].

\textbf{Compute budget.} Roughly 50 GPU hours for the main evaluation sweep, plus 20 GPU hours for sensitivity studies.

\begin{thebibliography}{99}
\bibitem{vllm} W. Kwon, Z. Li, S. Zhuang, Y. Sheng, L. Zheng, C. H. Yu, J. E. Gonzalez, H. Zhang, and I. Stoica. Efficient Memory Management for Large Language Model Serving with PagedAttention. \emph{SOSP}, 2023.
\bibitem{distserve} Y. Zhong, S. Liu, J. Chen, J. Hu, Y. Zhu, X. Liu, X. Jin, and H. Zhang. DistServe: Disaggregating Prefill and Decoding for Goodput-optimized Large Language Model Serving. \emph{OSDI}, 2024.
\bibitem{sarathi} A. Agrawal, N. Kedia, A. Panwar, J. Mohan, N. Kwatra, B. S. Gulavani, A. Tumanov, and R. Ramjee. Taming Throughput-Latency Tradeoff in LLM Inference with Sarathi-Serve. \emph{OSDI}, 2024.
\bibitem{llumnix} B. Sun, Z. Huang, H. Zhao, W. Xiao, X. Zhang, Y. Li, and W. Lin. Llumnix: Dynamic Scheduling for Large Language Model Serving. \emph{OSDI}, 2024.
\bibitem{splitwise} P. Patel, E. Choukse, C. Zhang, A. Shah, \'I. Goiri, S. Maleki, and R. Bianchini. Splitwise: Efficient Generative LLM Inference Using Phase Splitting. \emph{ISCA}, 2024.
\bibitem{mooncake} R. Qin, Z. Li, W. He, J. Cui, H. Tang, F. Ren, T. Ma, S. Cai, Y. Zhang, M. Zhang, Y. Wu, W. Zheng, and X. Xu. Mooncake: A KVCache-centric Disaggregated Architecture for LLM Serving. \emph{FAST}, 2025.
\bibitem{qlm} A. Patke, D. Reddy, S. Jha, H. Qiu, C. Pinto, C. Narayanaswami, Z. Kalbarczyk, and R. Iyer. Queue Management for SLO-Oriented Large Language Model Serving. \emph{SoCC}, 2024.
\bibitem{scorpio} Y. Tang, T. Lan, X. Huang, H. Lu, and W. Chen. Scorpio: Serving the Right Requests at the Right Time for Heterogeneous SLOs in LLM Inference. arXiv:2505.23022, 2025.
\bibitem{jitserve} W. Zhang, Z. Wu, Y. Mu, R. Ning, B. Liu, N. Sarda, M. Lee, and F. Lai. JITServe: SLO-aware LLM Serving with Imprecise Request Information. arXiv:2504.20068, 2025.
\bibitem{slosserve} S. Chen, Z. Jia, S. Khan, A. Krishnamurthy, and P. B. Gibbons. SLOs-Serve: Optimized Serving of Multi-SLO LLMs. arXiv:2504.08784, 2025.
\bibitem{banaserve} Y. He, M. Xu, J. Wu, J. Hu, C. Ma, M. Shen, L. Chen, C. Xu, L. Qu, and K. Ye. BanaServe: Unified KV Cache and Dynamic Module Migration for Balancing Disaggregated LLM Serving in AI Infrastructure. arXiv:2510.13223, 2025.
\bibitem{longbench} Y. Bai, X. Lv, J. Zhang, H. Lyu, J. Tang, Z. Huang, Z. Du, X. Liu, A. Zeng, L. Hou, Y. Dong, J. Tang, and J. Li. LongBench: A Bilingual, Multitask Benchmark for Long Context Understanding. \emph{ACL}, 2024.
\end{thebibliography}

\end{document}
